\section{The experimental setup}
\label{sec:setup}

We evaluated Qwen2.5-3B-Instruct, a 36-layer model, in fp32 with greedy decoding. The prompt asks it to write
a short memory note and then an \texttt{Answer:} line containing a city or ``unknown''. The instruction
explicitly offers ``unknown'' when the memory does not supply an answer.

\paragraph{What the note tells us.}
The note is a prompted summary of the memory, not a verified reasoning trace. When we describe a response as
``note, then refuse'', we mean that the gold city appears in that summary before \texttt{Answer: unknown}.
This shows that the city was written down. It does not establish that the model completed an internal
derivation, and we do not treat the note as a faithful explanation of the computation.

\paragraph{Tasks and examples.}
The synthetic memories support three tasks: \textbf{direct extraction (DE)}, such as ``Where does A live?'';
\textbf{relational composition (RC)}, such as ``In which city does A's supervisor work?''; and
\textbf{temporal updating (TU)}, such as ``Where does A live now?'' after a move.
Each underlying \emph{base case} produces 13 related rows: 6 answerable rows (two variants at three hardness
levels), 3 unanswerable rows with a related distractor (UNA-rel), 1 unanswerable row with a swapped name
(UNA-swap), and 3 counterfactual rows in which the evidence specifies a different city (CF).

We call the unedited model \textbf{B0}. Its responses define the populations used throughout the evaluation:
\textbf{FA} contains answerable rows on which B0 abstains falsely; \textbf{C} contains answerable rows it gets
correct; \textbf{UNA} contains all unanswerable rows; and \textbf{CF} contains the changed-evidence rows.
Each intervention is evaluated on the same fixed populations.

\paragraph{Splits.}
Related variants stay together when cases are assigned to exploration, selection, confirmation and transfer
splits (\cref{tab:splits}). The partition also prevents a person--evidence-city key, including a
counterfactual variant's key, from crossing splits. No canonical base key or exact memory text appears in two
splits. Names, cities and templates can still be shared; disjoint cases do not imply entirely different
vocabularies. The agent saw only the exploration split, E.

\begin{table}[ht]
\centering\small
\begin{tabularx}{\linewidth}{@{}llYY@{}}
\toprule
Split & Bases per task & Access & Purpose \\
\midrule
E & 48 & Agent & Exploration only \\
S & 32 & Analysts & Dose selection and candidate selection \\
H & 48 & Sealed; read once & Confirmation \\
T-city & 24 & Sealed; read once & Transfer to new cities \\
\bottomrule
\end{tabularx}
\caption{E and S use generator seed 51 with disjoint cases, H uses seed 52, and T-city uses seed 53.
Rows from earlier experiments are excluded by prompt hash. A second planned transfer split with a reworded
instruction, T-format, was dropped before the first design freeze (F0).}
\label{tab:splits}
\end{table}

\paragraph{Scoring.}
A frozen rule-based grader scores the answer after the last \texttt{Answer:} marker, separately from the note
(\cref{app:methods}). It distinguishes a correct answer, abstention, wrong answer, mixed answer and abstention,
conflicting answers, other content, and degenerate or missing outputs.

The primary utility, $U$, rewards a correct answer with $+1$ and assigns abstention $0$. A mixed answer receives
$0$ on an answerable row and $-1$ on an unanswerable row; other errors receive $-1$. Within each base,
answerable and unanswerable rows each contribute half the weight. Counterfactual rows are scored separately
and do not enter $U$. We also report:
\begin{itemize}
  \item \textbf{Recovery (Rel):} the share of baseline false abstentions that become correct answers.
  \item \textbf{Unsupported answering (Uns):} the share of unanswerable rows given an asserted or suggested answer.
  \item \textbf{Damage (Dam):} the share of baseline-correct rows that are no longer correct.
  \item \textbf{Counterfactual following (CF-follow):} the share of changed-evidence rows answered with the new city.
  \item \textbf{Degenerate outputs (Deg):} degenerate responses and responses missing an answer marker.
\end{itemize}
Rates are base-weighted. We also give raw row counts where useful; these need not yield the same percentages.

\paragraph{Comparators.}
Five interventions serve as repair baselines, with settings frozen after selection on S:
\begin{itemize}
  \item \textbf{B1, logit bias:} penalise ``unknown'' tokens after the answer marker.
  \item \textbf{B2, added instruction:} append ``Do not reply unknown.'' to the prompt.
  \item \textbf{B3, option removal:} remove the ``unknown'' option from the instruction.
  \item \textbf{B4, activation steering:} apply the earlier refusal-direction push at layer 29.
  \item \textbf{B5, directional orthogonalisation:} remove a direction fitted on E from the relevant weights,
        following an adaptation of \citet{arditi2024refusal}.
\end{itemize}
A sixth comparator, \textbf{B6}, consists of five size-matched random edits. It is used in the pathway tests,
not in the main repair comparison (\cref{sec:explain}).

\begin{figure}[ht]
\centering
\begin{tabularx}{\linewidth}{@{}YYYY@{}}
\toprule
\textbf{Explore} & \textbf{Select} & \textbf{Freeze} & \textbf{Evaluate} \\
\midrule
Agent proposes edits on E; the harness replays them. &
S selects candidates and doses separately for repair and pathway tests. &
F2 fixes candidates, settings and analysis before confirmation. &
H tests the frozen claims; T-city checks transfer descriptively. \\
\bottomrule
\end{tabularx}
\caption{The agent sees only E. S selects the candidates and doses for repair and pathway tests. The
confirmatory plan is frozen by commit before any model outputs on H or T are generated; each is then read once.}
\label{fig:pipeline}
\end{figure}
